\section{Introduction}
UAVs inspect power lines, map crops, deliver medicine and extend cellular coverage after disasters
\cite{zeng2016mag,mozaffari2019}. Most of these uses make an existing service faster or cheaper. UAVs also enable
operations that did not exist before because no infrastructure could support them: urgent, one-off missions in places
where the network is weak.

Searching for a missing person is such a mission. A lost hiker or a missing child has often walked past many cameras
installed by different owners at different times. A few of them, or an edge box attached to them, can run a neural
matcher; most are legacy devices that only record. None knows whom to look for. Where an AI camera has broadband, the
query can simply be sent to it; our concern is the large remainder in mountain resorts, parks and along rural roads, where
devices are linked only by low-power radios and footage cannot be uploaded. We therefore push a compact query down
instead of pulling footage up. In a network built by many owners the AI nodes run different embedding models whose
vectors are not comparable \cite{shen2020}; the query must carry reference images so that each node embeds them with its
own model. We use $L=12$ images, about 400\,kB, i.e., $Q=4041$ packets, as the nominal case.

\textbf{Challenge.} A fixed-wing UAV covers a large area but cannot hover; it crosses each cluster once. One copy of the
query takes about 2\,km of flight. If the air-to-ground (A2G) range is about 1\,km (open land), nodes far from the path
still receive many packets, but scattered, so few complete an image. If the range is 100--200\,m (canopy, tall buildings),
a cell hears only the few hundred packets emitted while the UAV is overhead. Either way the AI nodes, which are rarely on
the path, cannot assemble the query alone. The classical answers are to make the UAV dwell near every receiver until it
has decoded a rateless code \cite{zeng2018multicast,luby2002}, or to let the ground disseminate the object epidemically
\cite{hui2004,levis2004}. The first costs airtime at every receiver; the second discovers at run time what the base
station (BS) could have known in advance.

\textbf{Key insight.} Packet $s$ is emitted at a point $\gamma(d_0+s\Delta)$ of a path $\gamma$ that the BS designs.
Before take-off the BS therefore knows which segment of the stream lands on which cell, and can plan both the path and the
ground cooperation: \emph{the flight geometry is the transmission schedule}.

\textbf{Contributions.} This paper extends our conference version \todo{cite ICC 2027 submission} with a formal problem
statement, fundamental limits, a repair mechanism for trajectory errors and a substantially larger evaluation.
\begin{enumerate}
\item \emph{Formulation and complexity.} We state single-pass query delivery as a joint choice of a curvature-constrained
trajectory and a ground forwarding schedule, and prove that it is NP-hard already without ground links
(Section~\ref{sec:problem}).
\item \emph{Limits of non-cooperative delivery.} For every scheme in which a camera uses only what its own cell hears,
we prove that the completion time is at least $(1-\delta)\,Q\,\Tp\,\Lambda^\star$, where $\Lambda^\star$ is a fractional
covering number of the critical cells under the reception footprint, obtained from a linear programme. The bound holds
for any trajectory, any code and correlated fading, and yields a sufficient condition under which ground cooperation
is strictly faster, including the case of no-fly zones (Section~\ref{sec:limits}).
\item \emph{River.} A pre-planned, cell-level transport, \spine{}\,$\to$\,Bank\,$\to$\,Ring\,$\to$\,critical cell, that
needs no route or content discovery, together with \emph{takeover}, a repair that restores delivery when the real
trajectory deviates from the plan, a characterisation of when it is triggered, and a single-radio lower bound against
which we measure River's ground time (Section~\ref{sec:river}).
\item \emph{Flight design.} \chord{} and \lmtour{}, which choose the path for the ground transport rather than for
coverage, the latter guided by a surrogate of River's completion time learned from ns-3 (Section~\ref{sec:flight}).
\item \emph{Evaluation} in ns-3 with a full 802.15.4 MAC over three maps, four A2G regimes and up to ten network
realisations per setting: mission time, the measured schemes against the proven limits, frames and radio energy
against Deluge and planned direct pulling, and robustness to trajectory error (Section~\ref{sec:eval}).
\end{enumerate}

\section{Related Work}\label{sec:related}
\textbf{UAV-enabled multicasting.} UAV trajectories are routinely optimised jointly with communication
\cite{zeng2016mag,mozaffari2019,wu2018,zeng2017energy}. Closest to our baseline, Zeng \emph{et al.}
\cite{zeng2018multicast} minimise the time for every ground terminal to recover a common file, coded with random linear
network coding, directly from a rotary-wing UAV. They lower-bound the recovery probability by the time the UAV spends
within a distance $D$ of each terminal, which turns the design into a waypoint problem, and they name a subsequent
device-to-device sharing phase as future work. Sharing among users after an aerial multicast has since been studied for
one-hop D2D links in cellular settings \cite{wang2020d2d,miao2020}. We consider the opposite end of the design space: a
fixed-wing UAV that cannot dwell, receivers that are a known subset of a multi-hop 802.15.4 mesh, and cooperation that is
planned from the trajectory before take-off. Our lower bound (Theorem~\ref{thm:cover}) applies to the whole class of
direct-reception designs, including \cite{zeng2018multicast}.

\textbf{Bulk dissemination in sensor networks.} Deluge \cite{hui2004} and Trickle \cite{levis2004} spread an object to
every node by advertise--request epidemics. MNP lets nodes sleep while neighbours send segments they do not need
\cite{kulkarni2005}; Typhoon uses unicast with snooping and channel switching \cite{liang2008}; Splash and Pando build on
constructive interference and reach an order of magnitude lower dissemination time than Deluge on testbeds
\cite{doddavenkatappa2013,du2017}; the Low-Power Wireless Bus schedules all communication centrally over synchronous
floods \cite{ferrari2012}. All of them start from a sink and target every node. River differs in where the data enters
(a known segment at every cell under the path), in whom it targets (a few critical cells) and in that the complete
forwarding plan follows from the flight geometry. Fountain and network codes \cite{luby2002,shokrollahi2006,ho2006}
relax \emph{which} packets are needed but not \emph{where} they are; River is compatible with both. For a single
broadcast channel, Stojanovic \emph{et al.} \cite{stojanovic2009} compare coded retransmission with cooperation among
receivers and find a crossover in the erasure probability; our crossover is geometric and involves the trajectory.

\textbf{Curvature-constrained routing.} The Dubins traveling salesman problem (DTSP) is NP-hard \cite{leny2012}; with
neighbourhoods (DTSPN) it is solved in practice by sampling entry configurations and transforming the resulting
generalised TSP into an asymmetric TSP \cite{isaacs2011,noon1993,obermeyer2009}, and data-collection variants add travel
budgets \cite{penicka2019}. Our tours use these building blocks, but the objective is the readiness time of the ground
transport rather than tour length.

\textbf{A2G channels and flight accuracy.} Low-altitude A2G propagation depends strongly on environment and elevation
angle \cite{khawaja2019,alhourani2014,matolak2017}, and field measurements of 2.4\,GHz 802.15.4 links from UAVs report
ranges well below nominal \cite{nekrasov2019}. We therefore treat the range as a regime parameter. Small fixed-wing
aircraft follow planned paths within metres under moderate wind \cite{park2007,nelson2007}; turns planned in still air
are distorted by wind \cite{techy2009}, which motivates the repair of Section~\ref{sec:takeover}.

\textbf{Person search at the edge.} Face recognition and person re-identification compare a probe with stored crops by
cosine similarity of learned embeddings \cite{schroff2015,deng2019,ye2022}; embeddings of different models are
incompatible \cite{shen2020}. We use these results only to define the query and the value of partial data.
